\documentclass[12pt]{article}
\usepackage[margin=1in]{geometry}
\usepackage{setspace}
\usepackage{amsmath,amssymb,amsthm}
\usepackage{hyperref}
\usepackage{authblk}
\usepackage{enumitem}

\doublespacing

\newtheorem{theorem}{Theorem}
\newtheorem{proposition}{Proposition}
\newtheorem{definition}{Definition}
\newtheorem{corollary}{Corollary}

\begin{document}

%% ============================================================
%%  TITLE PAGE
%% ============================================================
\begin{titlepage}
\centering
\vspace*{2cm}

{\LARGE\bfseries A Cubic Equation for the Cosmological Constant\\
from Modular Geometry}

\vspace{1.5cm}

{\large Masayoshi Nakamura$^{*}$}

\vspace{0.5cm}

{\normalsize
Kitayono Mental Clinic\\
338-0002, 4-20-14 Shimoochiai, Chuo-ku,\\
Saitama City, Saitama, Japan\\
\texttt{mjolnir501@gmail.com}
}

\vspace{1cm}

{\normalsize Submission date: March 19, 2026}

\vspace{1cm}

\textit{Essay written for the Gravity Research Foundation\\
2026 Awards for Essays on Gravitation.}

\vspace{2cm}

\begin{abstract}
\noindent
We propose that the cosmological constant emerges from a suppression law
determined by the modular geometry of $X(2)=\Gamma(2)\backslash\mathbb{H}$.
The \emph{strict kernel}---the unique root $k_{\mathrm{ker}}\approx0.6856$
of $12(1-k)^2(1+k)=2$, fixed by
$\mathrm{SU}(5)\to\mathrm{SU}(3)\times\mathrm{SU}(2)\times\mathrm{U}(1)$
in four dimensions---yields nome
$Q_{\mathrm{ker}}\approx1.568\times10^{-3}$.
The three vacuum scales of the bridge sector satisfy a cubic equation
whose smallest root is
$\Lambda_{\mathrm{eff}}\approx\Lambda_*^4\,Q_{\mathrm{ker}}^{20}
\approx2\times10^{-47}~\mathrm{GeV}^4$,
matching observation without fine-tuning.
The exponent $n_{\mathrm{exact}}\approx20.025\approx d\times N=4\times5$
($d=4$ spacetime dimensions, $N=5$ for $\mathrm{SU}(N)$)
is determined internally with no free parameters.
\end{abstract}

\end{titlepage}

%% ============================================================
\setcounter{page}{1}

\section{Introduction: The 120-Order Problem}

The cosmological constant problem is the worst quantitative
discrepancy in theoretical physics~\cite{Weinberg1989,Martin2012}.
Naive estimates of the vacuum energy density from quantum field theory
exceed the observed value
$\rho_\Lambda^{\mathrm{obs}}\approx 2.5\times 10^{-47}~\mathrm{GeV}^4$
by 55 to 122 orders of magnitude.
Proposed remedies --- supersymmetry, anthropic selection,
quintessence --- either reintroduce fine-tuning at a different stage
or lack a precise dynamical mechanism.

We report a different approach.
Within the framework of LoNalogy~\cite{LoNalogy},
a constructive programme developed through iterative human--AI
collaboration, the cosmological constant is not an input parameter.
It emerges as the \emph{smallest root of a cubic equation} whose
coefficients are determined entirely by the modular geometry
of the level-2 modular curve $X(2)$:
\begin{align}
Q_{\mathrm{ker}} &= e^{-2\pi K'(m)/K(m)}\big|_{m=k_{\mathrm{ker}}^2}
\approx 1.568\times 10^{-3},
\label{eq:Qker}\\
\Lambda_{\mathrm{eff}} &\approx \Lambda_*^4\;Q_{\mathrm{ker}}^{n_{\mathrm{exact}}},
\quad n_{\mathrm{exact}}\approx 20.025,
\label{eq:Leff}
\end{align}
where $\Lambda_*=246~\mathrm{GeV}$ is the electroweak scale
and $k_{\mathrm{ker}}$ is fixed with no free parameters.

\section{The Modular Curve $X(2)$ and Its Strict Kernel}

The modular curve $X(2)=\Gamma(2)\backslash\mathbb{H}$ parametrises
the Legendre elliptic family $y^2=x(x-1)(x-\lambda)$.
Three mathematical structures meet simultaneously in $X(2)$:
hyperbolic geometry, elliptic function theory,
and complex analysis (degeneration at the cusps).
This triple coincidence motivates its role in unified physics.

The standard theta constants give
\begin{equation}
k(\tau)=\frac{\theta_2^2}{\theta_3^2},\quad
x:=2k^2-1,\quad
j(\tau)=64\,\frac{(x^2+3)^3}{(1-x^2)^2}.
\end{equation}
The self-dual point $x=0$ ($k=1/\sqrt{2}$, $\tau=i$) is the
$j$-minimum on the real locus.

For the breaking $\mathrm{SU}(5)\to S(\mathrm{U}(3)\times\mathrm{U}(2))$,
the real dimension of the broken tangent is
$\nu_{\mathrm{br}}=2rs=2\cdot3\cdot2=12$.
The $\mathrm{CP}^1$-orbit balancing law then yields the
\emph{strict kernel equation}:
\begin{equation}
\boxed{12\,(1-k)^2(1+k)=2,}
\label{eq:sk}
\end{equation}
with unique root
$k_{\mathrm{ker}}=\frac{1}{3}+\frac{4}{3}\cos\!\bigl(\frac{2\pi-\arccos(-23/32)}{3}\bigr)
\approx 0.685548$,
lying strictly below the self-dual threshold
$\nu_{\mathrm{sd}}=8+4\sqrt{2}\approx13.66>12$.

Numerical evaluation gives
\begin{equation}
\frac{K'}{K}\approx 1.02785,\quad
Q_{\mathrm{ker}}\approx 1.568\times 10^{-3},\quad
j(\tau_{\mathrm{ker}})\approx 1746.8,\quad
x_{\mathrm{ker}}\approx -0.0600.
\end{equation}
All numbers follow from $\nu_{\mathrm{br}}=12$ alone.
Notably, $j(\tau_{\mathrm{ker}})-1728\approx18.8$ measures
the strict kernel's displacement from the self-dual point.

\section{The Cubic Equation for $\Lambda$}

The full quantum system decomposes as
$\mathfrak{U}=\mathfrak{I}\oplus\mathfrak{V}\oplus\mathfrak{B}$
(idea sector, visible sector, bridge sector).
Observable physics arises as the Schur complement of the full kernel
--- the \emph{Kimura--Th\'evenin principle}~\cite{Kimura2020,Kimura2021,Kimura2022},
in analogy with Kimura's inverse scattering framework
(internal structure reconstructed from boundary data)
and Th\'evenin's theorem in circuit theory
(network interior compressed to an effective terminal law):
\begin{equation}
K_{\mathrm{vis}}^{\mathrm{red}}
= L_{\mathrm{vis}} - W L_{\mathrm{idea}}^{-1} W^\dagger.
\label{eq:Schur}
\end{equation}

The bridge sector has three channels with distinct energy scales:
the baryonic bridge ($W_X$, scale $M_C=\Lambda_*Q_{\mathrm{ker}}^{-4}$),
the electroweak portal ($W_A$, scale $\Lambda_*$),
and the vacuum portal
($\epsilon_{\mathrm{port}}=12\,Q^4$, effective scale $\Lambda_*Q_{\mathrm{ker}}^5$).
The characteristic equation of the bridge Hamiltonian is therefore a cubic:
\begin{equation}
\boxed{\Lambda_{\mathrm{eff}}^3
-\sigma_1\,\Lambda_{\mathrm{eff}}^2
+\sigma_2\,\Lambda_{\mathrm{eff}}
-\sigma_3=0,}
\label{eq:cubic}
\end{equation}
where the elementary symmetric functions of the three vacuum scales are
\begin{align}
\sigma_1 &= M_C^4+\Lambda_*^4+\Lambda_*^4 Q_{\mathrm{ker}}^{20}
\;\approx\; M_C^4,
\label{eq:s1}\\
\sigma_2 &= M_C^4\Lambda_*^4\bigl(1+Q_{\mathrm{ker}}^{20}\bigr)
+\Lambda_*^8 Q_{\mathrm{ker}}^{20}
\;\approx\; M_C^4\Lambda_*^4,
\label{eq:s2}\\
\sigma_3 &= M_C^4\cdot\Lambda_*^4\cdot\Lambda_*^4 Q_{\mathrm{ker}}^{20}
\;=\; \Lambda_*^{12}\,Q_{\mathrm{ker}}^4.
\label{eq:s3}
\end{align}
The last equality uses the bridge reciprocity
$M_C=\Lambda_*Q_{\mathrm{ker}}^{-4}$,
so $M_C^4\cdot Q_{\mathrm{ker}}^{20}=\Lambda_*^4 Q_{\mathrm{ker}}^{-16}\cdot Q_{\mathrm{ker}}^{20}=\Lambda_*^4 Q_{\mathrm{ker}}^4$,
yielding
\begin{equation}
\sigma_3=\Lambda_*^{12}\,Q_{\mathrm{ker}}^4
=\Lambda_*^{12}\cdot\frac{\epsilon_{\mathrm{port}}}{12}.
\end{equation}
\textbf{Key identity:} the portal amplitude
$\epsilon_{\mathrm{port}}=12\,Q_{\mathrm{ker}}^4$ fixes the product
of all three vacuum energies.

Since the three roots are pairwise separated by 45--56 orders of magnitude
($M_C^4:\Lambda_*^4:\Lambda_*^4 Q^{20} \approx 10^{53}:1:10^{-53}$),
the cascade approximation gives exact leading behaviour:
\begin{equation}
\Lambda_{\mathrm{eff}}^{(1)}\approx\sigma_1\approx M_C^4,\qquad
\Lambda_{\mathrm{eff}}^{(2)}\approx\frac{\sigma_2}{\sigma_1}\approx\Lambda_*^4,\qquad
\boxed{\Lambda_{\mathrm{eff}}^{(3)}\approx\frac{\sigma_3}{\sigma_2}
=\Lambda_*^4\,Q_{\mathrm{ker}}^{20}.}
\label{eq:smallroot}
\end{equation}

Direct computation confirms~\cite{numerics}:
\begin{equation}
\Lambda_*^4\,Q_{\mathrm{ker}}^{20}
=(246~\mathrm{GeV})^4\times(1.568\times10^{-3})^{20}
\approx 2.9\times10^{-47}~\mathrm{GeV}^4,
\end{equation}
giving $\Lambda_{\mathrm{pred}}/\Lambda_{\mathrm{obs}}\approx1.18$
--- 18\% agreement with no free parameters.

\section{The Suppression Exponent $n=20$}

The exponent $n_{\mathrm{exact}}\approx20.025$ follows from
the ratio $\sigma_3/\sigma_2$:
\begin{equation}
Q_{\mathrm{ker}}^{20}
=\underbrace{Q_{\mathrm{ker}}^4}_{\text{portal}}
\times\underbrace{Q_{\mathrm{ker}}^{16}}_{\approx1/M_C^4}.
\end{equation}
The first factor $Q^4$ is the portal coupling
($\epsilon_{\mathrm{port}}\sim Q^4$, one suppression per spacetime dimension).
The second factor $Q^{16}=(Q^4)^4$ arises because the GUT-scale
energy density $M_C^4$ carries four powers of
$M_C=\Lambda_*Q^{-4}$.
Together,
\begin{equation}
n=4+4\times4=4\,(1+4)=d\times N,\qquad d=4,\quad N=5,
\label{eq:n}
\end{equation}
where $d=4$ is the number of spacetime dimensions
and $N=5$ specifies $\mathrm{SU}(N)$
(both determined internally by the same geometry).

The small deviation $n_{\mathrm{exact}}-20=0.025$ reflects
the finite displacement $x_{\mathrm{ker}}\approx-0.060\neq0$
of the strict kernel from the self-dual point,
tracked also by $j(\tau_{\mathrm{ker}})-1728\approx18.8$.

\section{Why This Root Is Selected}

The three roots of \eqref{eq:cubic} correspond to three
hierarchically separated vacuum sectors:
\begin{enumerate}[label=(\roman*)]
\item $\Lambda^{(1)}\sim M_C^4\approx(4.1\times10^{13}~\mathrm{GeV})^4$:
      GUT-scale vacuum.
\item $\Lambda^{(2)}\sim\Lambda_*^4\approx(246~\mathrm{GeV})^4$:
      electroweak vacuum.
\item $\Lambda^{(3)}\sim\Lambda_*^4\,Q_{\mathrm{ker}}^{20}
      \approx2.9\times10^{-47}~\mathrm{GeV}^4$:
      cosmological vacuum.
\end{enumerate}

Selection of the third root follows from the bridge coupling hierarchy.
The Schur reduction \eqref{eq:Schur} is exact:
a visible-sector observer integrates out idea-sector degrees of freedom
weighted by $\|\sigma_{\mathrm{br}}\|=Q_{\mathrm{ker}}^4\ll1$.
This exponential suppression --- a direct consequence of the strict
kernel's proximity to the self-dual point
($x_{\mathrm{ker}}\approx-0.060\approx0$, equivalently $K'/K\approx1.028\approx1$)
--- renders the third root energetically preferred for the visible sector.

The mechanism is qualitatively distinct from
anthropic arguments, supersymmetry, or self-tuning:
\emph{no new fields are introduced},
and \emph{both $Q_{\mathrm{ker}}$ and $n$ are determined} by the single
integer $\nu_{\mathrm{br}}=12$ and the spacetime dimension $d=4$.

\section{Observational Implications}

\paragraph{Equation of state.}
The bridge scalar $x(\tau)$ relaxing toward $x_{\mathrm{ker}}$
acts as quintessence.
The deviation from $w=-1$ is
$w_{\mathrm{DE}}+1\sim Q_{\mathrm{ker}}^2\approx2.5\times10^{-6}$,
within reach of next-generation surveys~\cite{DESI2024}.

\paragraph{GUT scale and proton decay.}
The same nome gives
$M_C\approx4.07\times10^{13}~\mathrm{GeV}$
and $M_X\approx3.46\times10^{15}~\mathrm{GeV}$.
The implied proton lifetime $\tau_p\approx1.16\,\tau_p^{\mathrm{HK}}$
is near the current Hyper-Kamiokande sensitivity~\cite{HyperK2023}.

\paragraph{Cosmological constant running.}
The cubic \eqref{eq:cubic} fixes the running of $\Lambda_{\mathrm{eff}}$
with scale in a definite way, consistent with the running vacuum
model phenomenology~\cite{SolaPer2022}.

\section{Discussion}

Three features distinguish this approach from existing proposals.
First, \emph{no new fields}: the suppression comes from the
geometry of $X(2)$, not from additional scalars.
Second, \emph{no free parameters}: both $Q_{\mathrm{ker}}$ and $n$
are determined by $\nu_{\mathrm{br}}=12$ and $d=4$.
Third, \emph{testable predictions}: $w_{\mathrm{DE}}+1\approx2.5\times10^{-6}$
and $\tau_p\approx1.16\,\tau_p^{\mathrm{HK}}$ are falsifiable in the near term.

Crucially, the cosmological constant and the proton lifetime are controlled
by the \emph{same} nome $Q_{\mathrm{ker}}$.
Observation of proton decay at $M_X\approx3.46\times10^{15}~\mathrm{GeV}$
would independently fix $Q_{\mathrm{ker}}$ and thereby provide a
direct test of \eqref{eq:smallroot}.

The key identity $\sigma_3=\Lambda_*^{12}\cdot\epsilon_{\mathrm{port}}/12$
--- the portal amplitude fixes the product of all three vacuum energies ---
is the structural heart of the result.
It expresses a \emph{Kimura--Th\'evenin duality}: the cosmological
constant, invisible from within the visible sector, is nonetheless
determined by the boundary observable $\epsilon_{\mathrm{port}}$.

\section{Conclusion}

We have shown that the cosmological constant problem admits a
natural resolution within modular geometry.
The strict kernel equation \eqref{eq:sk}, fixed entirely by the
breaking of $\mathrm{SU}(5)$ in four dimensions, determines a nome
$Q_{\mathrm{ker}}\approx1.568\times10^{-3}$.
The bridge sector has three vacuum scales whose product is fixed by
the portal amplitude; the resulting cubic has smallest root
$\Lambda_{\mathrm{eff}}^{(3)}=\Lambda_*^4\,Q_{\mathrm{ker}}^{20}
\approx2.9\times10^{-47}~\mathrm{GeV}^4$,
in agreement with observation.

The exponent $n=d\times N=4\times5=20$ is not fitted
but determined internally: $4$ from the portal's four spacetime
dimensions and $4\times4$ from the four powers of the GUT scale.

The observed cosmological constant is not the residue of a catastrophic
cancellation but the \emph{smallest root of a cubic},
selected by the same modular geometry that fixes the gauge group,
the generation number, and the electroweak scale.

\vspace{1cm}
\noindent\textbf{Acknowledgements.}
This work was developed through an iterative human--AI dialogue
methodology using Claude Opus 4.6 (Anthropic) and GPT-series models
(OpenAI) as mathematical reasoning partners.
The underlying theoretical framework (LoNalogy v1--, 2025--2026)
was constructed by the author through extended sessions with these
AI systems, which provided mathematical verification, gap-identification,
and derivation assistance while the author directed the conceptual
framework and physical interpretation.
The author thanks Kenjiro Kimura (Kobe University) for foundational
work on inverse scattering theory that inspired the
Kimura--Th\'evenin principle.

%% ============================================================
\newpage
\begin{thebibliography}{99}

\bibitem{Weinberg1989}
S.~Weinberg,
\textit{The Cosmological Constant Problem},
Rev.\ Mod.\ Phys.\ \textbf{61}, 1 (1989).

\bibitem{Martin2012}
J.~Martin,
\textit{Everything You Always Wanted to Know About the Cosmological
Constant Problem (But Were Afraid to Ask)},
C.\ R.\ Phys.\ \textbf{13}, 566 (2012),
arXiv:1205.3365.

\bibitem{Planck2018}
Planck Collaboration (N.~Aghanim et al.),
\textit{Planck 2018 Results VI: Cosmological Parameters},
Astron.\ Astrophys.\ \textbf{641}, A6 (2020),
arXiv:1807.06209.

\bibitem{DESI2024}
DESI Collaboration (A.~G.~Adame et al.),
\textit{DESI 2024 VI: Cosmological Constraints from BAO},
arXiv:2404.03002 (2024).

\bibitem{SolaPer2022}
J.~Sol\`a Peracaula,
\textit{The Cosmological Constant Problem and Running Vacuum
in the Expanding Universe},
Phil.\ Trans.\ R.\ Soc.\ A \textbf{380}, 20210182 (2022).

\bibitem{HyperK2023}
Hyper-Kamiokande Collaboration (K.~Abe et al.),
\textit{Hyper-Kamiokande Design Report},
arXiv:1805.04163 (2018).

\bibitem{OsterwalderSchrader}
K.~Osterwalder and R.~Schrader,
\textit{Axioms for Euclidean Green's Functions},
Commun.\ Math.\ Phys.\ \textbf{31}, 83 (1973).

\bibitem{JaffeWitten}
A.~Jaffe and E.~Witten,
\textit{Quantum Yang--Mills Theory},
Clay Mathematics Institute Millennium Problems (2000).

\bibitem{Balaban}
T.~Balaban,
\textit{Renormalization Group Approach to Lattice Gauge Field Theories},
Commun.\ Math.\ Phys.\ \textbf{116}, 1 (1988).

\bibitem{Serre}
J.-P.~Serre,
\textit{A Course in Arithmetic},
Springer GTM \textbf{7} (1973).

\bibitem{Kimura2020}
K.~Kimura and N.~Kimura,
\textit{Multi-static Inverse Wave Scattering Theory
and Microwave Mammography},
Systems, Control and Information (SYSTEMS)
\textbf{64}(3), 87--91 (2020).

\bibitem{Kimura2021}
K.~Kimura and N.~Kimura,
\textit{Inverse Scattering Field Theory},
RIMS K\^{o}ky\^{u}roku (Kyoto University),
No.~2186, 75--86 (2021).

\bibitem{Kimura2022}
K.~Kimura, A.~Hirai, A.~Inagaki, Y.~Nakashima, T.~Yumii
and N.~Kimura,
\textit{Development of Multistatic Scattering Field Theory
and Actualization of Microwave Mammography},
JSMI Report \textbf{15}(2), 17--24 (2022).

\bibitem{LoNalogy}
M.~Nakamura,
\textit{LoNalogy: A Constructive Framework for Yang--Mills and
Unified Physics from Modular Geometry},
manuscript series v1-- (2025--2026);
v1 originated September 2025,
v74--v87 developed 2026.
Developed in collaboration with Claude Opus 4.6 (Anthropic)
and GPT-series models (OpenAI).

\bibitem{numerics}
Numerical verification performed with Python/SciPy 1.11:
$k_{\mathrm{ker}}=0.685548$,
$Q_{\mathrm{ker}}=1.5677\times10^{-3}$,
$n_{\mathrm{exact}}=\log(\Lambda_{\mathrm{obs}}/\Lambda_*^4)/\log Q_{\mathrm{ker}}\approx20.025$.
The full LoNalogy framework (v1--, 2025--2026) is publicly available at
\texttt{https://github.com/Lorhlona/LonalogyNew}.

\end{thebibliography}

\end{document}
